Nullniveau der Lageenergie: der tiefste Punkt. Zustand 1: im höchsten Punkt (nur Lageenergie). Zustand 2: beim Auftreffen auf die Matte (Lageenergie und kinetische Energie). Zustand 3: im tiefsten Punkt (nur Spannenergie).

\begin{abcliste}
\abc
Sei $h$ die Fallhöhe bis zur Matte und $v$ ihre Geschwindigkeit dort. Zwischen den Zuständen 1 und 2 fällt sie um $h$:
\begin{align}
 m\,g\,h &= \frac{1}{2}\,m\,v^2 \\
 v &= \vF \\
   &= \sqrt{2\cdot\g\cdot\hF} \\
   &= \v \approx \result{\vP{0}{2}}
\end{align}

\abc
Sei $s$, wie tief die Matte eingedrückt wird, und $k$ ihre Federkonstante. Von Zustand 1 bis Zustand 3 fällt sie um $h + s$ (während die Matte eingedrückt wird, sinkt sie weiter), und die ganze Energie wird zur Spannenergie der Matte:
\begin{align}
 m\,g\,(h+s) &= \frac{1}{2}\,k\,s^2 \\
 \frac{1}{2}\,k\,s^2 - m\,g\,s - m\,g\,h &= 0
\end{align}
Die positive Lösung dieser quadratischen Gleichung:
\begin{align}
 s &= \sMF \\
   &= \frac{\m\cdot\g+\sqrt{\left(\m\cdot\g\right)^2+2\cdot\k\cdot\m\cdot\g\cdot\hF}}{\k} \\
   &= \sM \approx \result{\sMP{0}{2}}
\end{align}
\end{abcliste}
